The source explicitly defines $v_p$ as the factor coprime to $p$, rather than the usual exponent. For $p=2$, $1+\sum_{i=0}^{100}2^i=2^{101}$, whose coprime part is 1. For each odd prime, the expression is 2 modulo $p$, so its coprime part is the whole expression $1+\dfrac{p^{101}-1}{p-1}$. Therefore the exact requested sum is $\boxed{1+\sum_{p\in\{3,5,7,11,13,17,19,23,29,31,37,41,43,47,53,59,61,67,71\}}(1+\dfrac{p^{101}-1}{p-1})}$. Its decimal expansion is supplied in the computation audit. If the intended definition were instead the conventional exponent of $p$, the answer would be 101, contributed only by $p=2$.
